% Pursuing Stacks (cotes 134-2 to 134-8): archive cote 134-4, site page 35 (PDF page 36), typescript page 123.
% typescript. Draft compared: ps3.tex:1734-1810.
% AI-transcribed by Claude Opus 5.5 (claude-opus-5-5, adaptive thinking, effort high), Message Batches API, batch ps2 (msgbatch_01NVwr1Yuo8qJf5rykeGBXPK), 2026-11-12; not yet checked by eye.
\scanpage{134-4/35}
\mspage{123}
considerably weakened,  but these will be sufficient):

a)  Finite products exist in M\hand{. (``pour mémoire'').}

b)  Arbitrary sums exist in M,  they are ``disjoint'' and ``universal''
(which implies that M has a \uline{strict} initial o\over{b}{b}ject).

c)  E\over{v}{v}ery object in M is isomorphic to a direct sum of O-connected one\hand{s.}

I recall an object is called \uline{O-connected} if it is  a)``non empty'' i.e.
non isomorphic to $\emptyset_M$ ,  and b) connected,  i.e. any decomposition of it
into a sum of two \inserted{sub}objects  is trivial (namely, one is ``empty'',  the other
is ``full''). Also,  under the assumption b),  well known standart argu-%
ments show that for any object\inserted{X}, a decomposition of X into a direct sum
of \over{O}{\unsure{0}}-connected subobjects is essentially uniqqe (namely the \inserted[typed]{corresponding} \unsure{\deleted{set of}|set of} subob-%
jects of X\inserted{is unique}), if it exists\over{,}{.} \deleted{is necessarily unique.} The  subobjects \deleted{that}
occuring in the sum are called the \uline{connected components} of X ,  and the
set of connected components of X is denoted as usual by $\hand{\pi}_o$(X).  We thus
get a natural functor
\[ X \hand{\mapsto} \xout{xx(X)}\,\hand{\pi}_o(X) \quad , \qquad M \hand{\to} (Sets) \quad \hand{.} \]
\xout{xxxxx}\inserted[typed]{This} can be equally described as a \uline{left} \uline{adjoint} to the functor
\[ E \hand{\longmapsto} E_M \quad , \qquad (Sets) \hand{\longrightarrow} M \quad , \]
associating to a set E the corresponding ``\uline{constant} \uline{object}'' E$_M$ of M,  some-%
times also denoted by the product symbol Exe$_M$ = Exe  (e the final object
of M ),  namely the direct sum \inserted{in \hand{M}} of E copies of e.  The adjunction formula
is
\[ \mathrm{Hom}_M(X,E_M) \hand{\simeq} \mathrm{Hom}_{(Sets)}(\hand{\pi}_o(X),E\ ) \quad . \]
The adjunction \xout{formula} relation implies that  the functor $\hand{\pi}_o$  commutes
with \inserted{all} direct limits  which exist in M,  and in particular (and trivially so)
to direct sums.

I'll finally assume also,  to fix the ideas :

d)  The final object of M is O-connected i.e. $\hand{\pi}_o$(e) = one point set.

Empty objects of M,\xout{xxx} on the other hand,  are of course characterized
by the condition
\[ \hand{\pi}_o(\emptyset_M) = \emptyset \quad . \]

Finally,  I'll make in the end a very strong assumption on M , which
however is satisfied more or less trivially in the cases we are interested
in,  when M is a would-be modelizer:

e)  (Total O-asphericity of M):  the product of two O-connected ob-%
jects of M is again O-connected.

This is clearly equivalent to the condition

e')  The functor $\hand{\pi}_o$ : M $\hand{\to}$(Sets)  commutes to finite products.

\margin{35}
